\cotefile{34}
\scanpage{180}
\mspage{15}
% archive sheet number: 179
% block: top, line 1 | text
3.2. Démonstration de \unsure{1.4|1.9} à l'aide de 3.1.

% block: upper third, lines 2-7 | text
Comme $U_0$ est \unsure{l'ouvert des|l'un des} $t_0\in B_0$ \unsure{dont} la fibre $\overline X_{t_0}$ est lisse, on sait que $f$ applique $U$ dans $U_0$; donc, \deleted{\unsure{l'inertie est}} \deleted{$\unsure{\pi_1}$} en vertu du th. du changement de base pour les $\pi_1^{(p)'}$ (SGA 1 \unsure{XIII}\ldots) on est ramené à prouver le th., \deleted{\unsure{dans}} dans la situation \unsure{avec indices 0}. Or, en vertu d'une \inserted{variante du} lemme de Abhyankar (SGA 1 XIII\ldots), le groupe fondamental de $U_0$ est \deleted{abélien premier} abélien d'ordre premier à $p$.

% block: middle third, lines 8-11 | text
Cela achève donc la démonstration de 1.4 dans le cas \deleted{De façon précise, il est isomorphe} de la dim relative 1, \inserted{de 1.4 et de 1.3} donc dans le cas général, grâce aux réductions du n° 2.

% block: middle to lower third, lines 12-18 | text
\uline{Remarque} 3.2.1. De façon précise, la variante indiquée du lemme de Abhyankar affirme que l'on trouve un revêtement universel de $U$ en prenant \inserted{la limite} des revêtements de $U$ induits par les revêtements kummeriens $S[T_1,\ldots,T_{\unsure{r|p}}]/(T_1^m,\ldots,T_{\unsure{r|p}}^m)$, \uline{avec $m$} \uline{premier à} $p$. Donc on trouve un isom canonique

% block: lower third, line 19 | display
\[
\pi_1(U_0,\unsure{\zeta|\xi})\simeq \deleted{\unsure{\mathbb Z}}\mathbb Z'_{\deleted{\unsure{\ell}}}(1)^I,
\]

% block: lower third, lines 20-22 | text
$I$ étant l'ens d'indices pour \deleted{\illegible{1}} les pts \deleted{\unsure{singuliers}}\inserted{de non} lissité de $X_0$, et $\mathbb Z'_{\deleted{\unsure{\ell}}}(1)$ est le produit des $\mathbb Z_\ell(1)$ ($\ell$ nb premier, \inserted{avec} $\ell\ne p$), $\leftarrow\mathbb Z_\ell(1)=\varprojlim_m\mu_{\ell^m}(k)$.

% block: bottom, lines 23-25 | text
On peut donc dire que l'action de $\pi_1(U_0)$ est déterminée essentiellement par \deleted{un}\inserted{la donnée} \deleted{\unsure{certain}} d'\unsure{automorphismes} \inserted{$h_i$} de $\pi_1^{p'}(X_{\overline{\unsure{K|\eta|s}}})$, \deleted{\unsure{définis à}}\inserted{\unsure{couplés|accouplés}} (i.e. le choix d'un
